\chapter{dsDNA Core Review Notes and Answer Guidance}

\section{Facilitating Without Doing the Project for the Student}
In dsDNA Core sessions, evidence stays at the center before answers. The strongest teaching question is often, ``Which table, row, script, or citation supports that?'' Students learn to diagnose problems by moving from output to source, not by guessing.

\section{Daily Check-In Structure}
This five-minute structure works at the start or end of each day:
\begin{enumerate}
  \item What did you produce yesterday?
  \item What evidence supports the main thing you think you learned?
  \item What is uncertain, broken, missing, or confusing?
  \item What will you produce by the end of today?
  \item What decision do you need from the program lead or dsDNA Core mentor?
\end{enumerate}

\section{Expected Answers by Program Stage}
\begin{longtable}{p{0.18\textwidth}p{0.36\textwidth}p{0.36\textwidth}}
\toprule
\textbf{Stage} & \textbf{Strong student answer} & \textbf{Needs coaching}\\
\midrule
\endfirsthead
\toprule
\textbf{Stage} & \textbf{Strong student answer} & \textbf{Needs coaching}\\
\midrule
\endhead
Question design & Names the system, comparison, chemical evidence, and limitation. & Says only that they want to see what chemicals are important.\\
Database interpretation & Distinguishes identity, source annotation, occurrence, pathway, hazard, and assay evidence. & Treats any database hit as equivalent support.\\
GC-MS workflow & Explains match factor, retention time, base peak, area, and tentative identity caveats. & Says the software identified the compound without caveat.\\
Enrichment & Starts with validation, source coverage, and provenance. & Jumps directly to the most exciting trait or pathway.\\
Trait matrix & Explains source fields, confidence, matrix keys, and why weak traits were removed. & Uses every available column because more data seem better.\\
Figure interpretation & Ties each figure to a question and a source table. & Makes decorative plots without a decision or claim.\\
Literature integration & Uses papers to contextualize results and limitations. & Searches only for papers that support the desired conclusion.\\
Final report & Methods can be rerun and claims are traceable. & Reads like a narrative without reproducible evidence.\\
\bottomrule
\end{longtable}

\section{Common Program Responses}
\begin{tabularx}{\textwidth}{p{0.30\textwidth}Y}
\toprule
\textbf{Student pattern} & \textbf{Program response}\\
\midrule
Overly broad project & The work is limited to one chemical list, one comparison, and one main evidence family.\\
Too much time cleaning messy data & The student switches to a program-reviewed case dataset for training, then returns to the real data later.\\
Interpreting before validation & A validation summary screenshot or table is reviewed before biological meaning is discussed.\\
Matrix is too wide & The matrix is narrowed to medium-confidence traits tied to the research question.\\
Weak writing & The student writes one claim, one evidence sentence, and one limitation sentence for each figure.\\
No progress after errors & The workflow moves to a small reproducible example with the exact error, input, expected output, and actual output.\\
\bottomrule
\end{tabularx}

\section{Review Notes for Real-World Exercises}
\subsection{Unmatched Chemical}
A good answer does not delete unmatched chemicals silently. It should preserve them in a missingness table and identify the next action: synonym search, identifier correction, manual review, or follow-up validation.

\subsection{Hazard Code Trap}
A good answer separates hazard screening from risk interpretation. Students should state that risk requires exposure, dose, route, organism, and context.

\subsection{Pathway Does Not Mean Presence}
A good answer rewrites ``this pathway is active'' as ``this compound has KEGG pathway context consistent with...'' unless independent pathway activity evidence exists.

\subsection{Natural Product Occurrence}
A good answer separates direct taxonomic evidence from broad natural product context. Students should not imply occurrence in their organism unless the source supports it.

\subsection{Matrix Sparsity}
A good answer favors a smaller matrix if it has clearer interpretation, better confidence, and sufficient coverage across compounds. Students should explain what was removed and why.

\section{When To Slow the Program Down}
The program slows down when:
\begin{itemize}
  \item Students cannot explain their research question without reading notes.
  \item Source coverage is weak and students are ignoring it.
  \item Scripts contain absolute paths or manual steps that cannot be repeated.
  \item Figures are being made before variables are defined.
  \item Literature is being used as decoration rather than evidence.
\end{itemize}

\section{When To Accelerate}
The program accelerates when:
\begin{itemize}
  \item The student can rerun their workflow from a clean R session.
  \item The grouping dictionary is reviewed and matrix outputs are stable.
  \item Figures have clear captions and limitations.
  \item The student can explain what would change their interpretation.
\end{itemize}

\section{Final Review Questions}
Before the final package is accepted, the review asks:
\begin{enumerate}
  \item Which result is strongest, and what evidence supports it?
  \item Which result is most uncertain, and why?
  \item Which chemicals were excluded, unmatched, or weakly supported?
  \item Which database source was most useful, and which was least useful?
  \item What should the next student do first?
  \item What claim should not be made from this project?
\end{enumerate}
